% 集中定义 CFE 的 TikZ 样式和模块
\tikzset{
  % CFE连线样式
  cfe edge/.style={
    -{Latex[length=2mm]}, % 箭头样式
    draw=teal!50!black, % 连线颜色
    line width=.55pt % 连线宽度
  },
  % CFE模块外框样式
  cfe frame/.style={
    rounded corners=6pt, % 外框圆角
    draw=gray!45, % 外框边框色
    fill=white, % 外框背景色
    inner sep=3.2mm % 外框内边距，调小可减少CFE留白
  },
  % 定义 CFE 模块的样式
  pics/cfe module/.style={
    code={ % 定义可复用的 CFE 模块
      % 四个x坐标控制特征片之间的横向距离
      % 数值越接近，特征片越靠拢，越像堆叠效果
      \def\xin{0} % 输入片x坐标
      \def\xfone{1.35} % 256片x坐标
      \def\xftwo{2.45} % 128片x坐标
      \def\xfthree{3.35} % 64片x坐标
      %
      % \cfeCuboid参数：名字/x/正面宽/高度/厚度/标签/颜色
      % 高度和厚度决定片的大小，正面宽度决定片的薄厚感
      \cfeCuboid{input310}{\xin}{.1}{2.00}{1.5}{310}{yellow} % 输入310维方片
      \cfeCuboid{feat256}{\xfone}{.09}{1.60}{1.1}{256}{yellow} % 特征256维方片
      \cfeCuboid{feat128}{\xftwo}{.08}{1.22}{0.72}{128}{yellow} % 特征128维方片
      \cfeCuboid{feat64}{\xfthree}{.07}{.86}{0.42}{64}{yellow} % 特征64维方片
      %
      \node[font=\scriptsize, text=black] at (1.95, -.75) {}; % 模块标签
      %
      \begin{scope}[on background layer] % 只把外框画到背景层
        % fit列出所有特征片的边界点，外框会自动包住它们
        \node[cfe frame, fit=(input310-fit-a)(input310-fit-b)(feat256-fit-a)(feat256-fit-b)(feat128-fit-a)(feat128-fit-b)(feat64-fit-a)(feat64-fit-b)] (cfebox) {}; % CFE外框
      \end{scope}
      %
      \node[
        anchor=south west, % 左下角作为定位点
        font=\small, % 标题字号
        text=black % 标题颜色
      ] at ([xshift=2mm,yshift=1mm]cfebox.north west) {CFE Module}; % 模块标题
    }
  }
}
